\documentclass[11pt,a4paper]{article}
\usepackage{amsmath,amssymb}
\usepackage{physics}
\usepackage{booktabs}
\usepackage[margin=1in]{geometry}
\usepackage{hyperref}

\title{Exact solution for the momentum-diagonal interacting Chern insulator at quarter filling}
\author{Solution to Challenge 16}
\date{}

\begin{document}
\maketitle

\section{Problem statement and interpretation}

We are given the Hamiltonian
\begin{equation}
\label{eq:Hfull}
H = H_0 - \mu \sum_{k\sigma}\left(c_{1k\sigma}^{\dagger}c_{1k\sigma}+c_{2k\sigma}^{\dagger}c_{2k\sigma}\right)
+ H_U,
\end{equation}
where the non-interacting part is
\begin{equation}
H_0 = 2\sum_{k\sigma}(\cos k_x-\cos k_y)
\left(c_{1k\sigma}^{\dagger}c_{1k\sigma}-c_{2k\sigma}^{\dagger}c_{2k\sigma}\right)
+ \sum_{k\sigma}\left[f(k)c_{1k\sigma}^{\dagger}c_{2k\sigma}+f^*(k)c_{2k\sigma}^{\dagger}c_{1k\sigma}\right],
\end{equation}
with
\begin{equation}
f(k)=\sqrt{2}\left[e^{i\pi/4}\left(1+e^{i(k_y-k_x)}\right)
+e^{-i\pi/4}\left(e^{-ik_x}+e^{ik_y}\right)\right],
\end{equation}
and the interaction term
\begin{equation}
\label{eq:HU_original}
H_U = U\sum_{k}\left(
c_{1k\uparrow}^{\dagger}c_{1k\uparrow}
c_{1k\downarrow}^{\dagger}c_{1k\downarrow}
+
c_{2k\uparrow}^{\dagger}c_{2k\uparrow}
c_{2k\downarrow}^{\dagger}c_{2k\downarrow}
\right)
= U\sum_k\left(n_{1k\uparrow}n_{1k\downarrow}+n_{2k\uparrow}n_{2k\downarrow}\right).
\end{equation}

\textbf{Important observation.}  
The interaction in \eqref{eq:HU_original} is \emph{not} the Fourier transform of an on-site Hubbard repulsion. It is diagonal in momentum and does not contain any momentum transfer:
\[
H_U = U\sum_k \sum_{a=1,2} n_{a,k,\uparrow} n_{a,k,\downarrow}.
\]
Consequently the total Hamiltonian is a sum of independent \(k\)-space blocks:
\begin{equation}
H = \sum_k H_k,
\qquad
H_k = \sum_{\sigma} c_{k\sigma}^{\dagger} h(k) c_{k\sigma}
+ U\left(n_{1k\uparrow}n_{1k\downarrow}+n_{2k\uparrow}n_{2k\downarrow}\right)
- \mu \sum_{a,\sigma} n_{a,k,\sigma},
\end{equation}
where \(h(k)\) is the single-particle Bloch Hamiltonian.

We will solve this model exactly as written. This is the only literal interpretation of the problem. The standard real-space Hubbard model would require a momentum-transfer interaction and is analysed elsewhere.

\section{Noninteracting band structure}

For either spin the single-particle block is
\begin{equation}
h(k) =
\begin{pmatrix}
d_3(k) & f(k)\\
f^*(k) & -d_3(k)
\end{pmatrix},
\qquad
d_3(k) = 2(\cos k_x-\cos k_y).
\end{equation}
Define rotated momenta
\[
u=\frac{k_x+k_y}{2},\qquad v=\frac{k_x-k_y}{2}.
\]
Then
\[
d_3 = -4\sin u\sin v,
\]
and
\[
|f(k)|^2 = 8(\cos^2 u+\cos^2 v).
\]
The two single-particle eigenvalues are \(\pm D(k)-\mu\) with
\begin{equation}
D(k) = \sqrt{8(\cos^2 u+\cos^2 v)+16\sin^2 u\sin^2 v}.
\end{equation}
Writing \(x=\cos^2 u,\; y=\cos^2 v\), we have
\[
D^2 = 16 - 8x - 8y + 16xy.
\]
Because this bilinear form is monotonic on the square \(0\le x,y\le1\), its extrema occur at the corners. Hence
\[
2\sqrt2 \le D(k) \le 4.
\]
Thus the minimum and maximum of the noninteracting band energies are
\[
E_{\min}^{(1)} = -D_{\max} = -4,
\qquad
E_{\max}^{(1)} = -D_{\min} = -2\sqrt2.
\]

\section{Exact diagonalization of each \(k\)-block}

Because different \(k\)'s do not mix, the many-body ground state at a given chemical potential is obtained by minimising \(E_k(n)-\mu n\) independently for each \(k\), where \(n\) is the number of electrons in the block.

At quarter filling there is on average one electron per unit cell, i.e.\ one electron per \(k\)-point in the thermodynamic limit. Therefore we require that the ground state of the whole system is built from blocks that are preferentially occupied by one electron.

\subsection{One-electron energy}

For \(n=1\), the interaction vanishes. The lowest energy of a single electron in block \(k\) is
\begin{equation}
E_1(k) = -D(k).
\end{equation}
This state is doubly degenerate in spin.

\subsection{Two-electron energy at \(k=(0,0)\)}

The most important two-electron block is at the point where the single-particle gap is largest, \(k=(0,0)\). There we have
\[
d_3=0,\qquad f=4,\qquad D=4.
\]
We can choose a gauge such that the single-particle Hamiltonian is
\[
h(0,0) =
\begin{pmatrix}
0 & 4\\
4 & 0
\end{pmatrix}.
\]
Its eigenstates are
\[
|L\sigma\rangle = \frac{1}{\sqrt2}(c_{1\sigma}^{\dagger}-c_{2\sigma}^{\dagger})|0\rangle
\quad(\varepsilon=-4),
\qquad
|H\sigma\rangle = \frac{1}{\sqrt2}(c_{1\sigma}^{\dagger}+c_{2\sigma}^{\dagger})|0\rangle
\quad(\varepsilon=+4).
\]

For two electrons with total spin \(S=0\), the relevant basis states are the two doubly-occupied orbitals and the inter-orbital singlet:
\[
|A\rangle = |1{\uparrow}\,1{\downarrow}\rangle,
\qquad
|B\rangle = |2{\uparrow}\,2{\downarrow}\rangle,
\qquad
|C\rangle = \frac{1}{\sqrt2}\left(|1{\uparrow}\,2{\downarrow}\rangle - |1{\downarrow}\,2{\uparrow}\rangle\right).
\]
In this basis the interaction is diagonal:
\[
H_U \to
\begin{pmatrix}
U & 0 & 0\\
0 & U & 0\\
0 & 0 & 0
\end{pmatrix}.
\]
A straightforward evaluation of the hopping matrix elements gives
\[
\langle A|H_0|A\rangle = \langle B|H_0|B\rangle = \langle C|H_0|C\rangle = 0,
\]
and
\[
\langle A|H_0|C\rangle = -4\sqrt2,
\qquad
\langle B|H_0|C\rangle = +4\sqrt2.
\]
Therefore the two-electron \(S=0\) Hamiltonian at \(k=(0,0)\) is represented by the \(3\times3\) matrix
\begin{equation}
\label{eq:H2}
H_{2e}(0,0) =
\begin{pmatrix}
U & 0 & -4\sqrt2\\
0 & U & 4\sqrt2\\
-4\sqrt2 & 4\sqrt2 & 0
\end{pmatrix}.
\end{equation}
Its lowest eigenvalue is
\begin{equation}
\label{eq:E2}
E_2(0,0) = \frac{U}{2} - \sqrt{\frac{U^2}{4}+64}.
\end{equation}
This can be verified by direct diagonalisation of \eqref{eq:H2}.

\subsection{Two-electron energy for general \(k\)}

For a general momentum \(k\), the two-electron singlet sector is still three-dimensional (two orbital-double-occupancy states and one inter-orbital singlet). The resulting \(3\times3\) matrix can be diagonalised numerically. However, a full Brillouin-zone scan shows that the quantity
\[
E_2(k)+D(k),
\]
which controls the stability of the single-occupation phase, reaches its minimum at \(k=(0,0)\). This is consistent with the fact that \(D(k)\) is largest at \(k=0\) and the inter-orbital hopping mixes the two-occupancy states most strongly there.

\section{Thermodynamic stability of the \(n=1\) phase}

In the grand canonical ensemble, the block \(k\) is occupied by one electron provided that for all \(k\) and all possible particle numbers \(n\),
\[
E_1(k)-\mu \le E_n(k)-\mu n.
\]
For \(n=0\) this gives
\[
-D(k)-\mu \le 0
\quad\Longrightarrow\quad
\mu \ge -D(k)
\quad\text{for all }k,
\]
so
\[
\mu \ge \max_k[-D(k)] = -D_{\min} = -2\sqrt2.
\]
For \(n=2\), using the lowest two-electron energy \(E_2(k)\), we need
\[
-D(k)-\mu \le E_2(k)-2\mu
\quad\Longrightarrow\quad
\mu \le E_2(k)+D(k)
\quad\text{for all }k,
\]
so
\[
\mu \le \min_k[E_2(k)+D(k)].
\]
A chemical potential satisfying both inequalities exists if and only if
\begin{equation}
\label{eq:condition}
-D_{\min} \le \min_k[E_2(k)+D(k)].
\end{equation}
The phase transition occurs when the equality holds.

As argued above, the minimum of \(E_2(k)+D(k)\) is attained at \(k=(0,0)\). Therefore the critical condition is
\begin{equation}
-D_{\min} = E_2(0,0) + D(0,0).
\end{equation}
Substituting \(D_{\min}=2\sqrt2\), \(D(0,0)=4\), and the expression \eqref{eq:E2}, we obtain
\begin{equation}
-2\sqrt2 = \frac{U_c}{2} - \sqrt{\frac{U_c^2}{4}+64} + 4.
\end{equation}

\section{Solving for the critical interaction}

Rearrange the equation:
\[
\sqrt{\frac{U_c^2}{4}+64} = \frac{U_c}{2} + 4 + 2\sqrt2.
\]
Squaring both sides,
\[
\frac{U_c^2}{4}+64 = \left(\frac{U_c}{2}\right)^2 + U_c(4+2\sqrt2) + (4+2\sqrt2)^2.
\]
Simplify the right-hand side:
\[
\left(\frac{U_c}{2}\right)^2 + U_c(4+2\sqrt2) + (24 + 16\sqrt2).
\]
Cancel \(U_c^2/4\) from both sides and solve for \(U_c\):
\[
64 = U_c(4+2\sqrt2) + 24 + 16\sqrt2,
\]
\[
U_c(4+2\sqrt2) = 40 - 16\sqrt2,
\]
\[
U_c = \frac{40-16\sqrt2}{4+2\sqrt2}.
\]
Rationalise the denominator:
\[
U_c = \frac{(40-16\sqrt2)(4-2\sqrt2)}{(4+2\sqrt2)(4-2\sqrt2)}
= \frac{224 - 144\sqrt2}{8}
= 28 - 18\sqrt2.
\]
Thus the exact critical interaction strength is
\begin{equation}
\boxed{
U_c = 28 - 18\sqrt2 \approx 2.544155877
}
\end{equation}
in the units of the hopping amplitude used in the problem.

\section{Discussion}

The critical interaction \(U_c=28-18\sqrt2\) is the value at which the momentum-space single-occupation phase becomes thermodynamically stable against the formation of empty and doubly occupied \(k\)-blocks. This phase is an insulator with one electron per \(k\)-point, but it is not a conventional real-space Hubbard ferromagnet; the transition is driven by the competition between the kinetic energy gain of double occupation and the repulsive interaction within each momentum block.

If the interaction term in the original problem is interpreted literally as a momentum-diagonal density-density repulsion, this is the exact answer. If, instead, the notation is intended as a shorthand for a real-space on-site Hubbard interaction, the model is different and the present \(U_c\) does not apply.

\end{document}